\section{Step 11 Preserve the search across a nested call}\label{h263-motion-s01}
Begin with a frame sketch. Mark where you will retain the caller's return address,
the request values, the result pointer, the candidate coordinates and the current
best result. A candidate comparison reuses four argument registers, so the
original request pointer cannot remain live only in \texttt{\$a0}. The SAD entry
also has no obligation to preserve your scratch registers.

Write the prologue and matching epilogue before the loops. Check their saved
register addresses and keep the stack suitably aligned. Return through the saved
caller continuation, rather than through the continuation written by the last
nested call. Do not use an unallocated area below the stack pointer for locals.

\paragraph{Observe} A valid radius-zero request makes one nested comparison.
The function returns to its outer caller with the original saved registers and
stack pointer, even though \texttt{\$ra} changed during the comparison.
\Needspace{8\baselineskip}
\paragraph{Check} Where can a candidate coordinate safely remain across SAD?
\begin{itemize}
\item[A.] In any scratch register, because the caller used it first.
\item[B.] In a preserved register or an allocated stack local.
\item[C.] Only in an input pixel byte.
\end{itemize}
